% ── §1 Introduction ──
\section{Introduction}
\label{sec:intro}

The prediction of residue--residue contacts from multiple sequence alignments
(MSAs) has emerged as one of the most consequential applications of
evolutionary sequence analysis, underpinning modern protein structure
prediction pipelines~\cite{jumper2021} and enabling the determination of
protein folds from covariation signal alone~\cite{weigt2009}.  The dominant
paradigm, Direct Coupling Analysis (DCA), models the joint probability
distribution of amino-acid sequences through a Potts Hamiltonian whose
coupling parameters are inferred by inverting a covariance matrix computed
from deep alignments~\cite{morcos2011}.  While methods based on sparse inverse
covariance estimation~\cite{jones2012}, pseudo-likelihood maximisation~\cite{ekeberg2013},
and related formulations have achieved remarkable precision at the top-$L/5$
long-range contact level, the mathematical objects that these methods produce---%
a dense matrix of coupling constants indexed by residue pairs and amino-acid
identities---remain largely opaque to formal analysis.  In particular, no prior
work has established machine-checkable proofs that the encoding layer through
which biological sequences enter the statistical machinery preserves the
structural properties that the downstream inference requires.

This paper addresses that gap by constructing a formally verified bridge
between digital-logic design theory and protein coevolution analysis.  Our
central observation is that the Karnaugh map, a rectangular grid representation
of Boolean functions introduced by Karnaugh in 1953~\cite{karnaugh1953},
provides a natural embedding of aligned biological sequences when combined
with reflected binary Gray-code ordering of the amino-acid alphabet.  Under
this embedding, adjacent cells on the map correspond to pairs of residues that
differ in exactly one bit position of their Gray codes, which in turn
correspond to physicochemically similar amino acids under an appropriate
ordering of the twenty standard residues~\cite{he2012}.  The Quine--McCluskey
algorithm~\cite{quine1952,mccluskey1956} then extracts minimal sum-of-products
(SOP) expressions from thresholded frequency maps, producing sets of prime
implicants that we interpret as coevolutionary rules of the form ``if residue
$a$ appears at position $i$ and residue $b$ at position $j$, then positions
$i$ and $j$ are coevolutionarily constrained.''

The contribution of this work is threefold.  First, we formalise the entire
analytical chain in the Lean~4 theorem prover~\cite{lean4}, establishing
118 machine-checked theorems that cover Gray-code bijectivity and adjacency,
$k$-mer indexing injectivity, amino-acid encoding properties, contact-map
symmetry and irreflexivity, and implicant-cover soundness and completeness.
Second, we apply the verified pipeline to 150 protein domains from the PSICOV
benchmark~\cite{jones2012}, comprising over 1.6~million labelled residue-pair
observations against native experimental structures, and demonstrate that the
resulting Boolean circuits transfer across proteins: a universal contact
circuit trained on approximately one hundred proteins predicts native contacts
on held-out proteins at a mean top-$L/5$ precision of $0.179$, corresponding
to $5.1\times$ enrichment over random expectation.  Third, we show that the
Karnaugh-map cell-rate prior, when fused with mean-field DCA scores via simple
rank aggregation, yields a hybrid predictor that exceeds both individual
methods ($0.204$ vs.\ $0.169$ for DCA alone and $0.124$ for the circuit alone;
$p < 10^{-4}$ by Wilcoxon signed-rank test), providing the first demonstration
that Boolean-minimisation-derived features carry orthogonal structural
information not present in continuous coupling scores.

We further report honest negative results.  An attention-style ``tension''
mechanism, which normalises each position's mutual-information profile by its
strongest entropic partner, does not improve contact prediction beyond what
consensus identity features already provide.  Iterative direct information
computed from our mean-field implementation remains unreliable on families
with many fully conserved columns due to pseudo-inverse conditioning, and we
therefore adopt the APC-corrected Frobenius norm as the primary DCA score.
Finally, we assess reconstruction readiness against literature thresholds and
find that current circuit-level signal, while significantly above random
expectation, falls far short of the precision required for template-free
folding; the circuits are best understood as a compressed annotation layer
that identifies specific coevolutionary constraints rather than a replacement
for continuous physics-based scoring.

The paper is organised as follows.  Section~\ref{sec:theory-gray} introduces
the Gray-code encoding framework for nucleotide and amino-acid alphabets.
Section~\ref{sec:theory-kmap-qm} describes Karnaugh-map construction and
Quine--McCluskey minimisation with don't-care conditions.
Section~\ref{sec:theory-info} reviews the information-theoretic measures used
throughout.  Section~\ref{sec:theory-dca} summarises Direct Coupling Analysis.
Section~\ref{sec:theory-lean} presents the Lean~4 formal verification.
Sections~\ref{sec:methods-data}--\ref{sec:methods-validation} describe the
dataset, pipeline architecture, and validation framework.
Sections~\ref{sec:results-transfer}--\ref{sec:results-negative} present all
positive and negative results.  Section~\ref{sec:discussion} discusses findings
in context, and Section~\ref{sec:conclusion} concludes.
